% GENERATED FILE. Edit canonical lessons, then run npm run generate:books.
% canonical-source-hash: fnv1a64:83c00f6f26f9f290
% canonical-lessons: BN-C276-ghotona, BN-C276-kichhu, BN-C276-onek, BN-C276-manush, BN-C276-durghotona

\chapter{Something Happened}
\label{ch:bn-something-happened}
\hlchaptermodality{276}

\begin{chapteropening}
\textbf{By the end of this chapter:} \emph{I can say that something happened and how many people were there.}

Complete the last lesson of chapter 276: I can say that something happened and how many people were there.
\end{chapteropening}

\section[\texorpdfstring{\bn{ঘটনা} (ghôṭônā)}{ঘটনা (ghôṭônā)}]{\bn{ঘটনা} (ghôṭônā) — an event; an incident}

\label{lesson:BN-C276-ghotona}



\begin{quote}
Before the new one: say the Bengali for a female singer, then the Bengali for a female dancer.
\end{quote}

\subsection*{You'll want to know: \bn{ঘটনা}}
\textbf{\bn{ঘটনা}} — \emph{ghôṭônā} — \textquotedblleft{}an event; an incident\textquotedblright{}.

\begin{quote}
\bn{আমরা} \bn{কী} \bn{করব}, \bn{সেটা} \bn{ঘটনার} \bn{ওপর} \bn{নির্ভর} \bn{করে।}
\end{quote}

\begin{quote}
\emph{āmrā kī kôrbô, seṭā ghôṭônār opôr nirbhôr kôre.}
\end{quote}

\begin{quote}
What we do will depend on the incident.
\end{quote}

\bn{ঘটনা} \bn{ঘটে}, an event takes place: \bn{একটা} \bn{ঘটনা} \bn{ঘটল}, something happened.

\begin{cousinweb}
\textbf{\bn{ঘটনা}} (\emph{ghôṭônā}) is Sanskrit \emph{ghaṭanā}, taken over whole as a tatsama; Monier-Williams gives it senses from effort and motion to a taking effect. It sits beside the verb \bn{ঘটা}, to happen.
\end{cousinweb}

\begin{grammarlens}[title={What you've built}]
One more word for telling what happened.
\end{grammarlens}

\subsection*{Your turn}
\begin{itemize}
\raggedright
  \item \emph{Say it:} \emph{ghôṭônā}
  \item \emph{Say it:} \emph{ghôṭônā}, once more
  \item \emph{Say it:} \emph{āmrā kī kôrbô, seṭā ghôṭônār opôr nirbhôr kôre.} — What we do will depend on the incident.
  \item \emph{Recall:} say the Bengali for a female singer, then the Bengali for a female dancer, then say \emph{ghôṭônā} again
\end{itemize}

\begin{tcolorbox}[breakable,colback=teal!4,colframe=teal!35!black,title={Before you move on}]
What does \bn{ঘটনা} mean? (\textquotedblleft{}An event; an incident\textquotedblright{}.) Say it once more.
\end{tcolorbox}
\section[\texorpdfstring{\bn{কিছু} (kichhu)}{কিছু (kichhu)}]{\bn{কিছু} (kichhu) — something; some}

\label{lesson:BN-C276-kichhu}



\begin{quote}
Before the new one: say the Bengali for a female dancer, then the Bengali for an event.
\end{quote}

\subsection*{You'll want to know: \bn{কিছু}}
\textbf{\bn{কিছু}} — \emph{kichhu} — \textquotedblleft{}something; some\textquotedblright{}.

\begin{quote}
\bn{গতকাল} \bn{রাতে} \bn{পাড়ায়} \bn{কিছু} \bn{একটা} \bn{ঘটল।}
\end{quote}

\begin{quote}
\emph{gôtôkāl rāte pāṛāy kichhu ekṭā ghôṭlô.}
\end{quote}

\begin{quote}
Something happened in the neighbourhood last night.
\end{quote}

Before a noun it means some: \bn{কিছু} \bn{টাকা}, some money.

\begin{cousinweb}
\textbf{\bn{কিছু}} (\emph{kichhu}) comes down from Sanskrit \emph{kiṁcid}, 'anything' (CDIAL 3144, which lists Bengali \emph{kichu}). \bn{কিছু} \bn{একটা} is 'some one thing', something you cannot yet name.
\end{cousinweb}

\begin{grammarlens}[title={What you've built}]
One more word for telling what happened.
\end{grammarlens}

\subsection*{Your turn}
\begin{itemize}
\raggedright
  \item \emph{Say it:} \emph{kichhu}
  \item \emph{Say it:} \emph{kichhu}, once more
  \item \emph{Say it:} \emph{gôtôkāl rāte pāṛāy kichhu ekṭā ghôṭlô.} — Something happened in the neighbourhood last night.
  \item \emph{Recall:} say the Bengali for a female dancer, then the Bengali for an event, then say \emph{kichhu} again
\end{itemize}

\begin{tcolorbox}[breakable,colback=teal!4,colframe=teal!35!black,title={Before you move on}]
What does \bn{কিছু} mean? (\textquotedblleft{}Something; some\textquotedblright{}.) Say it once more.
\end{tcolorbox}
\section[\texorpdfstring{\bn{অনেক} (ônek)}{অনেক (ônek)}]{\bn{অনেক} (ônek) — many; a lot}

\label{lesson:BN-C276-onek}



\begin{quote}
Before the new one: say the Bengali for an event, then the Bengali for something.
\end{quote}

\subsection*{You'll want to know: \bn{অনেক}}
\textbf{\bn{অনেক}} — \emph{ônek} — \textquotedblleft{}many; a lot\textquotedblright{}.

\begin{quote}
\bn{ঘটনাটা} \bn{দেখতে} \bn{অনেক} \bn{লোক} \bn{রাস্তায়} \bn{এল।}
\end{quote}

\begin{quote}
\emph{ghôṭônāṭā dekhte ônek lok rāstāy elô.}
\end{quote}

\begin{quote}
A lot of people came out onto the road to see what had happened.
\end{quote}

\begin{cousinweb}
\textbf{\bn{অনেক}} (\emph{ônek}) is Sanskrit \emph{aneka}, \emph{an-} 'not' and \emph{eka} 'one' (Monier-Williams): 'not one' came to mean 'many'. \textbf{\bn{এক}}, one, is the same \emph{eka}.
\end{cousinweb}

\begin{grammarlens}[title={What you've built}]
One more word for telling what happened.
\end{grammarlens}

\subsection*{Your turn}
\begin{itemize}
\raggedright
  \item \emph{Say it:} \emph{ônek}
  \item \emph{Say it:} \emph{ônek}, once more
  \item \emph{Say it:} \emph{ghôṭônāṭā dekhte ônek lok rāstāy elô.} — A lot of people came out onto the road to see what had happened.
  \item \emph{Recall:} say the Bengali for an event, then the Bengali for something, then say \emph{ônek} again
\end{itemize}

\begin{tcolorbox}[breakable,colback=teal!4,colframe=teal!35!black,title={Before you move on}]
What does \bn{অনেক} mean? (\textquotedblleft{}Many; a lot\textquotedblright{}.) Say it once more.
\end{tcolorbox}
\section[\texorpdfstring{\bn{মানুষ} (mānush)}{মানুষ (mānush)}]{\bn{মানুষ} (mānush) — a human being; people}

\label{lesson:BN-C276-manush}



\begin{quote}
Before the new one: say the Bengali for something, then the Bengali for many.
\end{quote}

\subsection*{You'll want to know: \bn{মানুষ}}
\textbf{\bn{মানুষ}} — \emph{mānush} — \textquotedblleft{}a human being; people\textquotedblright{}.

\begin{quote}
\bn{ঘটনার} \bn{সময়} \bn{ওখানে} \bn{কিছু} \bn{মানুষ} \bn{ছিল।}
\end{quote}

\begin{quote}
\emph{ghôṭônār sômôy okhāne kichhu mānush chhilô.}
\end{quote}

\begin{quote}
There were a few people there when it happened.
\end{quote}

\begin{cousinweb}
\textbf{\bn{মানুষ}} (\emph{mānush}) comes down from Sanskrit \emph{mānuṣa}, 'human; a man' (CDIAL 10049, which lists Bengali \emph{mānus}). It names one person or people in general; \textbf{\bn{লোক}} is the plainer everyday word.
\end{cousinweb}

\begin{grammarlens}[title={What you've built}]
One more word for telling what happened.
\end{grammarlens}

\subsection*{Your turn}
\begin{itemize}
\raggedright
  \item \emph{Say it:} \emph{mānush}
  \item \emph{Say it:} \emph{mānush}, once more
  \item \emph{Say it:} \emph{ghôṭônār sômôy okhāne kichhu mānush chhilô.} — There were a few people there when it happened.
  \item \emph{Recall:} say the Bengali for something, then the Bengali for many, then say \emph{mānush} again
\end{itemize}

\begin{tcolorbox}[breakable,colback=teal!4,colframe=teal!35!black,title={Before you move on}]
What does \bn{মানুষ} mean? (\textquotedblleft{}A human being; people\textquotedblright{}.) Say it once more.
\end{tcolorbox}
\section[\texorpdfstring{\bn{দুর্ঘটনা} (durghôṭônā)}{দুর্ঘটনা (durghôṭônā)}]{\bn{দুর্ঘটনা} (durghôṭônā) — an accident}

\label{lesson:BN-C276-durghotona}



\begin{quote}
Before the new one: say the Bengali for many, then the Bengali for a human being.
\end{quote}

\subsection*{You'll want to know: \bn{দুর্ঘটনা}}
\textbf{\bn{দুর্ঘটনা}} — \emph{durghôṭônā} — \textquotedblleft{}an accident\textquotedblright{}.

\begin{quote}
\bn{গতকাল} \bn{সকালে} \bn{স্টেশনের} \bn{কাছে} \bn{একটা} \bn{বড়} \bn{দুর্ঘটনা} \bn{ঘটল।}
\end{quote}

\begin{quote}
\emph{gôtôkāl sôkāle sṭeshôner kāchhe ekṭā bôṛo durghôṭônā ghôṭlô.}
\end{quote}

\begin{quote}
There was a bad accident near the station yesterday morning.
\end{quote}

\begin{cousinweb}
\textbf{\bn{দুর্ঘটনা}} (\emph{durghôṭônā}) joins Sanskrit \emph{dur-}, 'bad, ill', to \textbf{\bn{ঘটনা}}, an event; Monier-Williams has both parts. A bad happening is an accident.
\end{cousinweb}

\begin{grammarlens}[title={What you've built}]
One more word for telling what happened.
\end{grammarlens}

\subsection*{Your turn}
\begin{itemize}
\raggedright
  \item \emph{Say it:} \emph{durghôṭônā}
  \item \emph{Say it:} \emph{durghôṭônā}, once more
  \item \emph{Say it:} \emph{gôtôkāl sôkāle sṭeshôner kāchhe ekṭā bôṛo durghôṭônā ghôṭlô.} — There was a bad accident near the station yesterday morning.
  \item \emph{Recall:} say the Bengali for many, then the Bengali for a human being, then say \emph{durghôṭônā} again
\end{itemize}

\begin{tcolorbox}[breakable,colback=teal!4,colframe=teal!35!black,title={Before you move on}]
What does \bn{দুর্ঘটনা} mean? (\textquotedblleft{}An accident\textquotedblright{}.) Say it once more.
\end{tcolorbox}
